%%%%%%%%%%%%%%%%%%%%%%%%%%%%%%%%%%%%%%%%%%%%%%%%%%%%%%%%%%%%%%%%%%
\section{Robustness}\label{sec:robustness}
%%%%%%%%%%%%%%%%%%%%%%%%%%%%%%%%%%%%%%%%%%%%%%%%%%%%%%%%%%%%%%%%%%
% Stream feeding this section: R3 (campaign/HANDOFF.md), and pre-campaign
% check B2 (collapsibleChannel solid resolution through AMI).

Partitioned Dirichlet--Neumann coupling converges more slowly, or not at all, as the added mass grows relative to the solid inertia, and Robin--Neumann coupling shifts but does not remove this limit~\citep{Causin2005, Forster2006, Tukovic2018b}.
The quasi-monolithic coupling solves the interface conditions within the Newton iteration, so its convergence should depend on the added mass only through the conditioning of the linear systems.
This section maps the outcome of each coupling over the density ratio and the time step on two cases with an exact or independent reference, so that a run that completes is also checked for accuracy.

\paragraph{Protocol.}
Each cell of a map is one run with one of four couplings: Aitken, IQN-ILS and Robin--Neumann under the strict convergence test with at most \tbd{100} coupling iterations per step, and the quasi-monolithic coupling.
For each run the outcome (Section~\ref{sec:protocol}), the mean number of Newton or coupling iterations per step, the wall time and the error of the quantity of interest against the reference are recorded.

%%--------------------------------------------------------------------------------------------------------------------%%
\subsection{Ring added-mass problem}\label{sec:robustness_ring}
%%--------------------------------------------------------------------------------------------------------------------%%
% Set-up as Section 4.2. Map: HANDOFF R3. Driver: ringAddedMass/verification/
% Allverify (density levels via its fluid-density setting); the extension
% beyond the strong level is new.

The ring of Section~\ref{sec:ring_m} is run on the factor-2 mesh with added-mass ratios from the weak to the strong level and beyond, down to the smallest solid-to-fluid density ratio at which any coupling completes, \tbd{range of $\mu$}, crossed with 12, 25, 50, 100 and 200 steps per period, for \tbd{number} periods.
The quantity of interest is the error of the wet/dry frequency ratio.

\todo{Figure: ringAddedMass robustness map. One panel per coupling (Aitken, IQN-ILS, Robin--Neumann, quasi-monolithic). Horizontal axis: steps per period (12, 25, 50, 100, 200; log). Vertical axis: added-mass ratio $\mu$ (log; weak, moderate, strong and the extension). Cell colour: outcome (completed, stalled, diverged, timed out); cell label: mean iterations per step; cells whose ratio error exceeds \tbd{threshold} hatched. Data: R3 \texttt{runs.csv}, \texttt{qois.csv}.}

Expected claims: the quasi-monolithic coupling completes every cell with \tbd{number} Newton iterations per step, independent of $\mu$, whereas the iterations of the partitioned couplings grow with $\mu$ and \tbd{which couplings} fail above \tbd{$\mu$} or below \tbd{steps per period}.

%%--------------------------------------------------------------------------------------------------------------------%%
\subsection{Collapsible channel}\label{sec:robustness_collapsible}
%%--------------------------------------------------------------------------------------------------------------------%%
% Set-up: companion paper sectionBenchmarkSuite.tex "Collapsible channel";
% solids4foam 257b2e564 tutorials/fluidSolidInteraction/collapsibleChannel/
% verification/README.md section "Quasi-monolithic coupling" (lines 293-335).
% Monolithic solid: nonLinearGeometryTotalLagrangianVelocity; conformal
% solid meshes 80x8 / 160x8 / 320x8 on fluid levels 1 / 2 / 4 at present.
% No monolithic driver: job scripts (A1 Section 1.4).

A thin, light elastic section of a channel wall collapses under a ramped external pressure and then oscillates~\citetodo{companion verification paper}; the problem modifies that of~\citep{Jensen2003} so that it is well posed for a continuum solid, and its reference is an independently converged oomph-lib solution.
It combines strong added mass with large deformation and fluid mesh motion, and the published problem has a massless wall, which Dirichlet--Neumann coupling cannot compute.

\todo{Figure: collapsible channel geometry (no figure exists in \texttt{figures/}; adapt the companion paper's or the tutorial's \texttt{images/} figure).}

\paragraph{Set-up.}
Fluid level 2 of the companion paper is used, with the wall computed by the geometrically nonlinear velocity solid on \tbd{a conformal $160\times8$ mesh / a finer solid mesh coupled through AMI with conservative traction transfer}.
The density ratio $\rho_s/\rho_f$ takes the values 10, 1, 0.1, 0.01 and 0 (massless wall), crossed with $\Delta t = 0.025$, $0.0125$, $0.00625$ and $0.003125\,\mathrm{s}$, and each run continues to the first trough of the wall midpoint or to $t = 1.5\,\mathrm{s}$.
The quantity of interest is the first trough of the wall-midpoint displacement, compared with oomph-lib at the same time step.
\todo{The quasi-monolithic tutorial variant uses Euler for the fluid \texttt{ddt(U)} (\texttt{system/fluid/fvSchemes.monolithic} at \texttt{257b2e564}); use \texttt{backward} (check the partitioned set-up of the companion paper), or the error against oomph-lib includes a first-order time error.}
\todo{Decision at gate G1 from check B2: with the linear-reconstruction velocity solid on the conformal solid meshes the trough error is set by the solid (Section~\ref{sec:lim_thin_solid}). If a finer solid mesh through AMI brings it to about 1\% at acceptable cost, report accuracy as well as robustness here and state the solid mesh; otherwise the map is a robustness result only and the hatching threshold is relaxed accordingly.}

\todo{Figure: collapsibleChannel robustness map, same layout as the ringAddedMass map. Horizontal axis: $\Delta t$ (0.003125 to 0.025 s; log). Vertical axis: $\rho_s/\rho_f$ (10, 1, 0.1, 0.01, 0; the massless row separated). Colour: outcome; label: mean iterations per step; hatching: trough error against oomph-lib above \tbd{threshold}.}

Expected claims: the quasi-monolithic coupling completes \tbd{which cells}, including the massless wall, with \tbd{iterations} per step, and the troughs for $\rho_s/\rho_f = 1$ and 0 agree to \tbd{value}; IQN-ILS fails \tbd{where} and Robin--Neumann for $\rho_s/\rho_f \le \tbd{value}$.
Where the quasi-monolithic coupling fails, the cause is \tbd{e.g.\ inversion of a fluid cell next to the wall by the predicted mesh displacement over one large step}, which is a limit of the mesh motion rather than of the coupled Newton iteration (Section~\ref{sec:limitations}).
% Recorded before the tag (stale; collapsibleChannel/verification/README.md,
% "Limitation"): monolithic level 2 failed at dt = 0.025 s at t = 0.25 s by
% fluid cell inversion. Report the campaign outcome, not this.

%%--------------------------------------------------------------------------------------------------------------------%%
\subsection{Flexible-bottom cavity: time step}\label{sec:robustness_cavity}
%%--------------------------------------------------------------------------------------------------------------------%%
% Optional in HANDOFF R3 (cheap). Level 3, dt = 0.0625 ... 1.0 s.

\todo{Optional (decide at G1). If run: cavityFlexibleBottom level 3 (Section~\ref{sec:cavity_m}) with $\Delta t = 0.0625$, $0.125$, $0.25$, $0.5$ and $1.0\,\mathrm{s}$, quasi-monolithic and Aitken. Here the time step is a pseudo-time step towards the steady state, and the measure is the number of steps and the wall time to the steady state as well as the outcome. Figure: wall time to steady state against $\Delta t$ for both couplings, with failed runs marked; or one row added to the maps above. Drop this subsection if the result adds nothing to Section~\ref{sec:robustness_collapsible}.}
